\subsection*{97. A 4-Week Learning Mechanism Experiment}

Consider a small, exploratory 4-week intervention for eligible engineers:

\begin{table*}[t]
\centering
\small
\caption{Structure of the 4-Week Learning Mechanism Experiment}
\begin{tabularx}{\textwidth}{p{0.18\textwidth}p{0.38\textwidth}X}
\toprule
\textbf{Stage} & \textbf{Learning Purpose} & \textbf{Example Evidence} \\
\midrule
Entry diagnostic & Locate misconceptions and current judgment & Unseen flawed evaluation case; stated reasoning \\
Week 1 & Build accurate causal models and baselines & Worked explanation, guided correction, independent variant \\
Week 2 & Practice contrasting failure cases & Mixed cases with feedback and required revision \\
Week 3 & Integrate in a production-style task & New logs, constraints, and a short design defense \\
Week 4 & Test independent application and transfer & Unseen variant, explicit assistance policy, rubric-based evaluation \\
\bottomrule
\end{tabularx}
\end{table*}

Important ideas should recur throughout the month. A separate delayed follow-up, where feasible and consented to, would help determine whether immediate gains persist. This is not a scientifically validated optimal duration. It is a bounded experiment that can be delivered and measured before the institution commits to a longer course.

\subsection*{98. A Teaching Session Is an Intervention, Not a Performance}

Imagine a live session where the founder shows 2 AI systems with similar apparent accuracy and asks which should be released. Learners commit to an individual answer, defend it, inspect new evidence, and revise. The instructor then explains the underlying evaluation principle, not simply the correct choice.

Next, each learner receives a different case that requires the same reasoning. The founder reads selected justifications, diagnoses common errors, and gives a revised task. A later session returns to the concept without warning that ``today is an evaluation lesson.''

The lecture's beauty is secondary. What matters is whether learners \textbf{notice the same consequential issue independently on a new task}. This distinction shapes faculty training, content production, review scheduling, and credible outcome claims.

\subsection*{99. Human Feedback Should Be Reserved for the Hardest Educational Judgment}

In a lean technical institute, instructor attention is scarce. Basic checks such as whether a script runs, whether required tests pass, or whether a simple probability calculation is correct can often be automated with well-tested solutions.

Human review should concentrate on causal reasoning, evaluation design, trade-off justification, and the ability to defend a decision. A structured rubric can focus the review: \emph{What was the learner's central hypothesis? Which evidence supports it? Which competing explanation was ignored? Which test would distinguish the alternatives?}

The institution should track how often feedback produces a meaningful revision. A 90-minute live class full of elegant explanations may be less useful than 30 minutes of well-targeted critique if the latter causes learners to repair a persistent mistake. Neither configuration should be assumed optimal without pilot evidence.

\subsection*{100. The Economics Must Include Feedback, Retries, and the Founder's Time}

Suppose an illustrative cohort has 12 learners. Each produces 4 submissions requiring 25 minutes of review, and half require a 15-minute 2nd review. Initial reviews require $12 \times 4 \times 25 / 60 = 20$ hours. Follow-ups add $24 \times 15 / 60 = 6$ hours, for \textbf{26 reviewer-hours} over the program. This excludes class time, preparation, questions, administration, and sales.

Those hours are not free simply because the founder does not initially draw a Bayes salary. For a mature model, the institution must be able to replace the work with sufficiently skilled paid faculty. Maintain both a \textbf{cash-cost ledger} for survival and a \textbf{normalized labor-cost ledger} for repeatability.

An apparently healthy course that depends on unlimited founder critique is not yet a scalable educational institution. A robust design reduces avoidable workload without destroying the diagnostic feedback that made the learning valuable.

\subsection*{101. The 1st Institute Should Resist Syllabus Inflation}

A founder with expertise in software, AI, mathematics, statistics, and quantitative finance can imagine an extraordinary 6-month curriculum. The danger is that every important subject becomes a compulsory module, crowding out the time learners need to practice, retrieve, revise, and transfer the fundamentals.

For each proposed topic, ask: \textbf{Which terminal performance requires it? What breaks if it is omitted? Can it be taught just in time? Is it a genuine prerequisite or an intellectually attractive extension?} A graduate who can defend a small set of difficult decisions may be better prepared than 1 who has heard a superficial survey of 40 technologies.

The richness of the institution's ambition should appear in \textbf{depth of reasoning and standards}, not only in breadth of its catalogue.

\subsection*{102. Selective Admissions Should Protect Learning, Not Manufacture Prestige}

If a course requires solid programming fundamentals, admission should establish those fundamentals. If weekly practice requires protected time, applicants should understand the time demands and show a plausible schedule. Selection can create a group able to work at a coherent level and support 1 another's learning.

But selectivity also changes evidence. If the strongest applicants perform exceptionally well at the end, that does not prove the instruction was exceptional. Bayes should retain baseline assessment data, report attrition honestly, and distinguish learners' prior ability from gains observed during the course.

A rigorous institute can tell an applicant ``not yet'' while offering a path to readiness. Refusal may protect the learner from a purchase that would waste money and effort; it need not be framed as a statement about their human potential.

\subsection*{103. The Institute's Brand Promise Creates an Educational Obligation}

``BECOME MORE.'' is a powerful aspiration. Its danger, if used carelessly, is that almost any course can claim to help people become something more. The mission becomes credible when it produces exact standards: \emph{Which capability, for which learner, in which context, proven how?}

An early Bayes promise should emphasize controllable educational obligations: honest prerequisites, competent instruction, meaningful attempts, timely feedback within stated limits, clear evaluation, and transparent results. It should not present a future campus, unearned institutional recognition, or employment guarantees as current product value.

The brand becomes a consequence of repeatedly delivering a rigorous transformation, not a substitute for doing so. A serious institution must be willing to show where it fell short and what it changed.

\subsection*{104. A Capability Transcript Could Be More Useful Than a Generic Certificate}

Over time, Bayes might develop a learner record that names specific competencies and links them to verified work. For each capability, the record might show entry state, demonstrated exit performance, tools allowed, whether an unseen task was used, and whether delayed retention was sampled.

This would require a calibrated assessment system, privacy safeguards, consent, and manageable faculty workload. An early pilot should not claim its assessment system already meets the evidentiary standards of an established examination board.

Still, the direction is important: \textbf{a credential should make the underlying competence more legible, not replace it with an uninspectable declaration}. A small institution can earn trust through transparent work samples and honest standards before its name alone is a recognized signal.
